\paragraph{Randomization of accelerated-branch assignments.}
To test whether the temporal distribution of branches alone could reproduce the observed positive $N_{\mathrm{eff}}^{+}$--diversification association, we generated 9,999 randomizations on the empirical 77-tip phylogeny. For each of 15 adult anatomical units, branches were classified as accelerated using the original inclusive 95th-percentile threshold separately for each of 1,000 branch-rate reconstruction samples. We independently permuted branch assignments within each unit, moving its complete 1,000-sample branch profile together. This preserved the number of accelerated branches in every sample and unit, while allowing assignments to move freely across all 152 branches. The phylogeny, branch durations and PyRate trajectory were held fixed; no branch-length matching was applied.

For every randomization and reconstruction sample, we recalculated $N_{\mathrm{eff}}^{+}=f^{+}\exp(H^{+})$ on the original 1,000-slice time grid, where $f^{+}$ is the proportion of intersecting branches with at least one accelerated unit and $H^{+}$ is the Shannon entropy of distinct nonzero whole-unit acceleration patterns among those branches. We averaged the resulting trajectories across reconstruction samples and then within thirteen complete 20-Myr bins from 260 to 0 Ma. Pearson correlations were calculated against the fixed family-level PyRate mean net-diversification rates. The upper-tail permutation $p$ value for the positive association was $(1+\#\{r_{\mathrm{shuffle}}\geq r_{\mathrm{observed}}\})/(9{,}999+1)$. Two-sided nominal Pearson $p$ values were also summarized, with the sign of the corresponding correlation retained. The figure shows all randomized trajectories and their arithmetic mean; no randomizations were selected or discarded.
